% Guidance for the Junior Level 1 show-me checklists and the week-1 chat
% (input by answer-key.tex; the pages the parent fills in are in assessments.tex)
\kcpart{A}{Assessments}{How to use the show-me checklists}

\begin{teachernote}[Watching, not testing]
\begin{itemize}[itemsep=1pt]
  \item \textbf{Nothing here is a grade.} The checklists tell you what to practice next and let you
        and your child see progress. Share them as \enquote{you can now\dots} and \enquote{next we will\dots}.
  \item \textbf{Watch during normal play.} Most rows show up on their own during the play day (week 5)
        or the project and showcase (weeks 9--10). Keep the sheet on a clipboard and check boxes as you see them.
  \item \textbf{Ask, do not quiz.} If a row has not come up, ask near the end:
        \enquote{Can you show me how you make the cat\dots?} Never time it, and never ask in front of
        a group if a child is shy: ask while you sit beside them.
  \item \textbf{Count what they do, not how they say it.} Pointing, acting it out and short answers all count.
        A child who reads slowly or does not write yet can still check \enquote{Got it}.
  \item \textbf{One box per row,} the best you saw that day. If you are unsure between two columns,
        choose the lower one and look again next week.
\end{itemize}
\end{teachernote}

%% one row of guidance: \smguide{row}{not yet}{getting there}{got it}{what to do next}
\newcommand{\smguide}[5]{%
  \par\needspace{7\baselineskip}\noindent\textbf{\color{kcaccent}#1}\par
  \begin{itemize}[leftmargin=1.4em,itemsep=1pt,topsep=2pt]
    \item[\faSeedling] \textbf{Not yet:} #2
    \item[\faShoePrints] \textbf{Getting there:} #3
    \item[\faCheckCircle] \textbf{Got it:} #4
    \item[\faArrowRight] \textbf{If not yet:} #5
  \end{itemize}\medskip}

%% ---------------------------------------------------------------------
\section{Week 1: \enquote{What do you like to make?} chat}
There are no right answers and no score. Use the chat like this:
\begin{itemize}[itemsep=2pt]
  \item \textbf{The circled pictures} (animals, games, stories, music, drawing) tell you which
        examples and free-play ideas to offer: an animal lover picks animal sprites, a music lover
        adds sounds first, a storyteller gets more \enquote{say} blocks.
  \item \textbf{The faces.} \enquote{Worried} or \enquote{not sure}: start them as Helper with a calm
        partner, praise effort out loud in weeks 1 and 2, and check in quietly after week 2.
        \enquote{Super excited}: give them a job such as the Checker, so the excitement turns into looking closely.
  \item \textbf{The drawing.} Mention it again in week 9 when they plan their greeting card.
\end{itemize}
Keep the page in your folder. It has no personal information beyond a nickname.

%% ---------------------------------------------------------------------
\section{Week 5 show-me checklist}

\smguide{Puts arrow cards in order to reach a goal (J1.G1)}
  {places cards at random, or walks the path without cards.}
  {gets most of the way, then fixes one card after trying it on the floor grid.}
  {lays out the whole path first, and the robot reaches the goal on the first or second try.}
  {Play the parent-robot game again with a short path (3 or 4 squares) and only forward and quarter-turn cards. Walk each card together as you lay it down.}

\smguide{Snaps a green-flag script together (J1.G2, J1.G3)}
  {drags blocks but they stay apart, or clicks single blocks instead of the flag.}
  {builds the script with reminders to start with the yellow hat block or to wait for the gray shadow.}
  {puts \enquote{when green flag clicked} on top, snaps 3 or more blocks under it and clicks the flag.}
  {Practice dragging: one block at a time, \enquote{wait for the shadow, then let go}. Build Hello Cat (week 2) together, with you as Helper.}

\smguide{Uses repeat to do something again (J1.G2)}
  {copies the same block several times instead of using repeat.}
  {uses repeat when reminded, or puts some blocks outside it by mistake.}
  {puts the steps inside a repeat, changes the number and predicts what will happen.}
  {Do the dance-loop game from week 4 again: \enquote{clap, stomp} four times with a repeat card around it. Then build a repeat with one block inside.}

\smguide{Finds and fixes a block that is not snapped on or is in the wrong order (J1.G4)}
  {says \enquote{it's broken} and waits for you to fix it.}
  {finds the bug when you point to the part of the script, then fixes it.}
  {watches what happens, finds the block, changes one thing and clicks the flag to test.}
  {Use the four steps on the cheat sheet out loud together: watch, find, change one thing, test. Put a working script beside the broken one to compare.}

\smguide{Takes turns as Driver and Helper (J1.G6, J1.G7)}
  {holds on to the mouse at the swap, or grabs it as Helper.}
  {swaps with a reminder; as Helper, sometimes reaches for the mouse.}
  {swaps when the timer runs out and, as Helper, points and uses words.}
  {Make the swap a ritual (say \enquote{Swap!}, change chairs). Give the Helper a job they like: reading the book, or checking blocks against the picture.}

\smguide{Says what their script does (J1.G6)}
  {shows the cat but cannot say what the blocks do yet.}
  {says what happens (\enquote{the cat talks}) but not which block does it.}
  {points to the blocks in order: \enquote{First it says hi, then it moves, then it meows.}}
  {In Share, ask one block at a time: \enquote{What does this one do?} Model it with your own script first.}

%% ---------------------------------------------------------------------
\section{Week 10 show-me checklist}

\smguide{Starts a script with the green flag, a click or a key (J1.G3)}
  {uses only the green flag, or puts a hat block in the middle of a script.}
  {uses a click or a key event when reminded where the yellow Events blocks are.}
  {chooses the right hat block for the job, for example a click to pop the balloon.}
  {Play \enquote{When I clap, you jump} from week 6, then add one \enquote{when this sprite clicked} script together.}

\smguide{Uses repeat for an animation or a dance (J1.G2)}
  {the sprite changes costume once, or the repeat is empty.}
  {the repeat works but goes too fast to see, or the number is not what they wanted.}
  {repeats \enquote{next costume} and a wait (or a turn) so the sprite dances, and changes the number on purpose.}
  {Show the week-4 bug again (too many or too few repeats). Add a wait inside the repeat together.}

\smguide{Uses two sprites or a backdrop change (J1.G2, J1.G5)}
  {has one sprite and one backdrop.}
  {adds a second sprite or a backdrop, but its script does not run.}
  {each sprite has its own script, or the backdrop changes when the card starts.}
  {Click each sprite below the stage and ask \enquote{What is in this sprite's code area?} Remind: each sprite has its own scripts.}

\smguide{Draws a shape with the pen (J1.G1, J1.G2)}
  {draws a line or nothing; pen down is missing.}
  {draws a square with help, or the turn is outside the repeat.}
  {uses pen down and repeat 4 of move and a quarter turn (90) to draw a square, or another shape from the turn card.}
  {Walk a square on the floor grid again: steps, quarter turn, four times. Then build it with the blocks beside the floor picture.}

\smguide{Plans with a storyboard and builds it (J1.G5)}
  {starts building without a plan, or the card does not match the drawing.}
  {has a storyboard and builds part of it, with reminders to look at it.}
  {draws Start, Middle and End, then builds one box at a time and checks it against the plan.}
  {Draw the three boxes together, one sentence each. Build only the Start box first and click the flag.}

\smguide{Finds and fixes a bug (J1.G4)}
  {asks you to fix it right away.}
  {finds the bug with a question such as \enquote{Which sprite is it in?}}
  {tests, finds and fixes it, and can say what was wrong.}
  {Use the four steps on the cheat sheet. Praise any step they do alone: \enquote{You watched very closely!}}

\smguide{Explains one script (J1.G6)}
  {says what the card does in general (\enquote{It's a birthday card}).}
  {names some blocks when asked about them.}
  {picks one script and says what each block does, in order, and why it is there.}
  {Practice the 1-minute show-and-tell script with you before the guests arrive.}

\smguide{Uses nicknames and gives credit (J1.G7)}
  {a real full name is on the card, or they say they made a library sprite.}
  {uses a nickname, and says where a picture came from when asked.}
  {uses a nickname or first name only and says, without being asked, \enquote{The balloon is from the Scratch library.}}
  {Talk about the week-6 and week-7 safety boxes again. Change the name on the card together.}

%% ---------------------------------------------------------------------
\section{What to do with the results}
\begin{itemize}[itemsep=2pt]
  \item \textbf{Week 5:} rows with \enquote{Not yet} for most children? Use buffer B1 (Unplugged Games
        Day) before week 6, or spend the next three warm-ups on those rows. One child only? Use the
        support path in weeks 6--8 and pair them as Helper with a patient partner.
  \item \textbf{Week 10:} keep the page with the certificate. Pass it on to the Junior Level 2 teacher
        (or keep it for yourself) for the week-1 warm-up of Level 2.
  \item \textbf{Celebrate growth:} put the week-5 and week-10 pages side by side with your child and
        count the boxes that moved to the right.
\end{itemize}
